% tikz settings
\usepackage{tikz}

\usetikzlibrary{
  shapes.geometric,
  arrows,
  positioning,
  fit,
  calc,
  shapes,
  backgrounds,
  matrix,
  pgfplots.groupplots
}
\usepackage{pgfplots}
\pgfplotsset{compat=1.18}
\usepackage{tikz-3dplot}
\tdplotsetmaincoords{70}{120}

\tikzset{
  block/.style={
      rectangle,
      rounded corners,
      minimum width=2.5cm,
      minimum height=1cm,
      text centered,
      draw=black,
      line width=1pt,
      fill=deepblue-background
    },
  circle/.style={
      draw=darkgray,
      line width=1pt,
      fill=darkgray-background,
      opacity=0.2
    },
  line/.style={thick, -, color=black},
  arrow/.style={thick, ->, >=stealth, color=black}
}

\pgfdeclarelayer{bg1}
\pgfdeclarelayer{bg2}
\pgfdeclarelayer{bg3}
\pgfdeclarelayer{fg1}
\pgfdeclarelayer{fg2}
\pgfdeclarelayer{fg3}
\pgfsetlayers{bg1,bg2,bg3,main,fg1,fg2,fg3}

% color settings
\definecolor{shallow-frame}{RGB}{211, 211, 211}
\definecolor{shallow-background}{RGB}{240, 240, 240}
\definecolor{warm-background}{RGB}{239, 239, 221}
\definecolor{blue-background}{RGB}{214, 220, 229}
\definecolor{deepblue-background}{RGB}{141, 170, 212}
\definecolor{darkgray-background}{RGB}{217, 217, 217}
\definecolor{mydarkgray}{RGB}{76, 76, 76}

% bibliography settings
% 适用于中文文章的 GBT-7714 参考文献格式
% https://ctan.org/pkg/gbt7714?lang=en
% https://github.com/zepinglee/gbt7714-bibtex-style
\usepackage{gbt7714}
\bibliographystyle{gbt7714-numerical}

\usepackage{booktabs}

% 章节标题设置
\ctexset{
  secnumdepth = 3,
  chapter = {
    format = \zihao{-2}\bfseries\raggedright,
    number = \chinese{chapter},
    name = {第,章}},
  section = {
    format = \zihao{-3}\bfseries},
  subsection = {
    format = \zihao{4}\bfseries}}

% 段落不加入编号
\renewcommand{\paragraph}[1]{
  \par\vspace{0.5\baselineskip}\noindent{\bfseries #1}\quad}

% 配置中文主字体：指定fonts目录下的Fandol字体
\usepackage{fontspec}

% 中文主字体：Fandol 宋体（含粗体、斜体配置）
\setCJKmainfont[
  Path       = fonts/,    % 字体文件所在相对路径（确保路径正确）
  Extension  = .otf,      % 字体文件扩展名
  BoldFont   = FandolHei-Regular,  % 粗体字体名
  ItalicFont = FandolKai-Regular   % 斜体（楷体）字体名
]{FandolSong-Regular}

% 中文无衬线字体：Fandol 黑体
\setCJKsansfont[
  Path       = fonts/,
  Extension  = .otf,
  BoldFont   = FandolHei-Bold
]{FandolHei-Regular}

% 中文等宽字体：Fandol 仿宋
\setCJKmonofont[
  Path       = fonts/,
  Extension  = .otf
]{FandolFang-Regular}

% 中文字族（同样修正参数格式）
\setCJKfamilyfont{zhsong}[
  Path       = fonts/,
  Extension  = .otf,
  BoldFont   = FandolSong-Bold
]{FandolSong-Regular}

\setCJKfamilyfont{zhhei}[
  Path       = fonts/,
  Extension  = .otf,
  BoldFont   = FandolHei-Bold
]{FandolHei-Regular}

\setCJKfamilyfont{zhfs}[
  Path       = fonts/,
  Extension  = .otf
]{FandolFang-Regular}

\setCJKfamilyfont{zhkai}[
  Path       = fonts/,
  Extension  = .otf
]{FandolKai-Regular}

% 字体命令
\renewcommand{\songti}{\CJKfamily{zhsong}}
\renewcommand{\heiti}{\CJKfamily{zhhei}}
\renewcommand{\fangsong}{\CJKfamily{zhfs}}
\renewcommand{\kaishu}{\CJKfamily{zhkai}}

\newenvironment{myproof}[1][\textbf{证明}]{\par\noindent{\heiti #1的证明}\quad}{\hfill$\square$\par\vspace{0.5\baselineskip}}